% ============================================================
%  MEDSTAND AI SALES ASSISTANT MANUAL — v3.0
%  Huong dan su dung + tai lieu chuc nang + tai lieu dao tao theo vai tro
%  Compile: XeLaTeX (2 lan)
% ============================================================

\documentclass[12pt,a4paper,oneside,openany]{book}

% --- Encoding & fonts ---------------------------------------------------------
\usepackage{fontspec}
\usepackage{polyglossia}
\setmainlanguage{vietnamese}
\setotherlanguage{english}
\setmainfont{Calibri}
\setsansfont{Calibri}
\setmonofont{Courier New}

% --- Layout ------------------------------------------------------------------
\usepackage[a4paper,top=2cm,bottom=2cm,left=2.4cm,right=2cm,headheight=14pt]{geometry}
\usepackage{parskip}
\setlength{\parindent}{0pt}
\setlength{\parskip}{4pt}

% --- Colors ------------------------------------------------------------------
\usepackage[dvipsnames,svgnames,table]{xcolor}
\definecolor{medgreen}{HTML}{0B8A43}
\definecolor{medblue}{HTML}{2563EB}
\definecolor{medindigo}{HTML}{4F46E5}
\definecolor{meddark}{HTML}{111827}
\definecolor{medgray}{HTML}{F3F4F6}
\definecolor{medline}{HTML}{D1D5DB}
\definecolor{medwarn}{HTML}{D97706}
\definecolor{mederr}{HTML}{DC2626}
\definecolor{roleTDV}{HTML}{0B8A43}
\definecolor{roleManager}{HTML}{D97706}
\definecolor{roleAdmin}{HTML}{2563EB}
\definecolor{roleCEO}{HTML}{7C3AED}
\definecolor{roleIT}{HTML}{374151}

% --- Graphics, tables, boxes -------------------------------------------------
\usepackage{graphicx}
\usepackage{tikz}
\usetikzlibrary{shapes,shadows,positioning,arrows.meta,calc}
\usepackage{booktabs}
\usepackage{longtable}
\usepackage{tabularx}
\usepackage{array}
\usepackage{multirow}
\usepackage[most]{tcolorbox}
\tcbuselibrary{skins,breakable}
\usepackage{fontawesome5}
\usepackage{amssymb}
\usepackage{enumitem}
\usepackage{titlesec}
\usepackage{fancyhdr}
\usepackage[colorlinks=true,linkcolor=medindigo,urlcolor=medindigo,citecolor=medgreen]{hyperref}
\usepackage{bookmark}
\usepackage{makeidx}
\makeindex

\titlespacing*{\chapter}{0pt}{0pt}{10pt}
\titlespacing*{\section}{0pt}{8pt}{4pt}
\titlespacing*{\subsection}{0pt}{6pt}{3pt}
\setlist[itemize]{leftmargin=2em,itemsep=2pt,topsep=2pt}
\setlist[enumerate]{leftmargin=2em,itemsep=2pt,topsep=2pt}

% --- Header & footer ---------------------------------------------------------
\pagestyle{fancy}
\fancyhf{}
\fancyhead[L]{\small\textcolor{medgreen}{\textbf{MEDSTAND}} — AI Sales Assistant Manual v3.0}
\fancyhead[R]{\small\leftmark}
\fancyfoot[C]{\small---\ \thepage\ ---}
\renewcommand{\headrulewidth}{0.4pt}

% --- Box styles --------------------------------------------------------------
\tcbset{
  noteBox/.style={enhanced,breakable,colback=blue!5,colframe=medblue,fonttitle=\bfseries\small,
    title={\faInfoCircle\ Ghi chú},attach boxed title to top left={yshift=-2mm,xshift=4mm},
    boxed title style={colback=medblue,rounded corners}},
  tipBox/.style={enhanced,breakable,colback=green!5,colframe=medgreen,fonttitle=\bfseries\small,
    title={\faLightbulb\ Mẹo vận hành},attach boxed title to top left={yshift=-2mm,xshift=4mm},
    boxed title style={colback=medgreen,rounded corners}},
  warnBox/.style={enhanced,breakable,colback=orange!8,colframe=medwarn,fonttitle=\bfseries\small,
    title={\faExclamationTriangle\ Lưu ý},attach boxed title to top left={yshift=-2mm,xshift=4mm},
    boxed title style={colback=medwarn,rounded corners}},
  exBox/.style={enhanced,breakable,colback=medindigo!5,colframe=medindigo!65,fonttitle=\bfseries\small,
    title={\faClipboardCheck\ Ví dụ thực tế},attach boxed title to top left={yshift=-2mm,xshift=4mm},
    boxed title style={colback=medindigo!75,rounded corners}},
  roleBox/.style={enhanced,breakable,colback=medgray,colframe=medline,left=7pt,right=7pt,top=7pt,bottom=7pt},
  codeBox/.style={enhanced,breakable,colback=meddark,colframe=meddark,fontupper=\ttfamily\small\color{green!35}}
}

% --- Helper commands ---------------------------------------------------------
\newcommand{\screenpurpose}[1]{\begin{tcolorbox}[colback=medgreen!8,colframe=medgreen!55,left=6pt,right=6pt,top=4pt,bottom=4pt]\faFlag\ \textbf{Mục đích:} #1\end{tcolorbox}}
\newcommand{\screenuser}[1]{\noindent\faUsers\ \textbf{Đối tượng sử dụng:} #1\par\smallskip}
\newcommand{\prompt}[1]{\begin{tcolorbox}[codeBox]#1\end{tcolorbox}}
\newcommand{\modulecard}[3]{\begin{tcolorbox}[roleBox]\textbf{#1}\par\smallskip\textbf{Người dùng:} #2\par\textbf{Giá trị:} #3\end{tcolorbox}}
\newcommand{\checkbox}{\makebox[1.2em][l]{$\square$}}

\begin{document}

% --- Cover -------------------------------------------------------------------
\begin{titlepage}
  \pagecolor{meddark}\color{white}\centering
  \vspace*{1.4cm}
  \begin{tikzpicture}
    \node[rectangle,rounded corners=16pt,fill=medgreen,minimum width=4.3cm,minimum height=4.3cm,
      drop shadow={shadow xshift=3pt,shadow yshift=-3pt,opacity=0.35}]{\Huge\faHospital};
  \end{tikzpicture}

  \vspace{0.8cm}
  {\fontsize{32}{40}\bfseries MEDSTAND AI SALES ASSISTANT}\\[0.35cm]
  {\Large CRM \& AI Copilot cho đội ngũ bán hàng ngành Dược}

  \vspace{0.9cm}
  \begin{tcolorbox}[colback=white!10!meddark,colframe=white!20!meddark,width=11cm,center]
    \centering\large\textbf{Tài liệu hướng dẫn sử dụng chính thức — Phiên bản 3.0}
  \end{tcolorbox}

  \vspace{0.7cm}
  \begin{tabular}{ccccc}
    \textcolor{roleTDV}{\faUser\ TDV} &
    \textcolor{roleManager}{\faUserTie\ Manager} &
    \textcolor{roleAdmin}{\faCogs\ Admin} &
    \textcolor{roleCEO}{\faChartLine\ CEO} &
    \textcolor{gray!70!white}{\faServer\ IT}
  \end{tabular}

  \vfill
  \begin{tcolorbox}[colback=medgreen!18!meddark,colframe=medgreen!45,width=10cm,center]
    \centering\small Tài liệu đào tạo nội bộ và triển khai khách hàng
  \end{tcolorbox}
  {\small\textcolor{gray!65!white}{\faCalendar\ Cập nhật: Tháng 7 năm 2026 \quad | \quad \faBuilding\ Medstand}}
\end{titlepage}
\pagecolor{white}\color{black}

% --- Document info -----------------------------------------------------------
\chapter*{Thông tin tài liệu}
\addcontentsline{toc}{chapter}{Thông tin tài liệu}

\begin{longtable}{p{4cm}p{10cm}}
\toprule
\textbf{Mục} & \textbf{Nội dung} \\
\midrule
Tên tài liệu & Hướng dẫn sử dụng MEDSTAND AI Sales Assistant \\
Phiên bản & 3.0 \\
Đối tượng & TDV/Sale, Quản lý vùng, Admin, CEO/Ban lãnh đạo, IT/System Admin \\
Mục đích & Đào tạo người dùng, mô tả chức năng, hỗ trợ vận hành và triển khai AI \\
Phạm vi & CRM, AI Copilot, tuyến bán hàng, đơn hàng, tồn kho, công nợ, CTBH, scoring, báo cáo \\
Mức độ bảo mật & Nội bộ / khách hàng triển khai \\
\bottomrule
\end{longtable}

\begin{tcolorbox}[noteBox]
Tài liệu này hợp nhất hai góc nhìn: \textbf{hướng dẫn thao tác hệ thống} cho người dùng cuối và \textbf{tài liệu chức năng} để BA, IT, quản lý dự án và đội triển khai hiểu đầy đủ phạm vi sản phẩm.
\end{tcolorbox}

\tableofcontents
\clearpage

% ============================================================================
\part{Tổng quan MEDSTAND AI}
% ============================================================================

\chapter{MEDSTAND là gì?}
\screenpurpose{Giúp người đọc hiểu MEDSTAND giải quyết bài toán nào trong hoạt động bán hàng ngành Dược.}

\textbf{MEDSTAND} là hệ thống CRM \& AI Sales Assistant dành cho ngành Dược, hỗ trợ đội ngũ trình dược viên và quản lý trong việc chăm sóc khách hàng, gợi ý đơn hàng, quản lý tuyến bán hàng, theo dõi tồn kho, tối ưu chương trình bán hàng và ra quyết định kinh doanh dựa trên dữ liệu.

Hệ thống gồm ba lớp chính:
\begin{enumerate}
  \item \textbf{Ứng dụng CRM/PWA}: dùng hằng ngày cho TDV, Manager và Admin.
  \item \textbf{AI Copilot}: trợ lý hỏi đáp, gợi ý bán hàng, phân tích tồn kho, công nợ, CTBH và dữ liệu khách hàng.
  \item \textbf{Lớp dữ liệu và tích hợp}: kết nối ERP, SQL Server, n8n, file PDF/Excel, catalog và ngân hàng tri thức nội bộ.
\end{enumerate}

\chapter{Mục tiêu triển khai AI}

\begin{longtable}{p{4cm}p{10cm}}
\toprule
\textbf{Mục tiêu} & \textbf{Ý nghĩa vận hành} \\
\midrule
Tăng doanh số & Gợi ý sản phẩm phù hợp theo từng khách hàng, lịch sử mua và chương trình bán hàng hiện tại. \\
Tăng hiệu suất TDV & Giúp TDV biết nên ghé ai, bán gì, khi nào ghé và cần nhắc khách vấn đề gì. \\
Cảnh báo tồn kho & Theo dõi tồn kho theo kho chính, DL02, DL03, lô hàng, hạn dùng và quyền truy cập của user. \\
Tối ưu tuyến & Ưu tiên khách sắp hết hàng, khách lâu chưa mua, khách có khả năng đạt thưởng hoặc có nguy cơ giảm mua. \\
Tối ưu CTBH & Nhắc khách còn thiếu bao nhiêu để đạt thưởng, chiết khấu hoặc mốc tích lũy. \\
\bottomrule
\end{longtable}

\begin{tcolorbox}[tipBox]
Khi demo cho khách hàng, nên bắt đầu bằng câu hỏi thực tế: ``Hôm nay TDV nên ghé khách nào và nên bán gì?'' Đây là tình huống thể hiện rõ nhất giá trị của AI.
\end{tcolorbox}

\chapter{Dữ liệu đầu vào của hệ thống AI}

AI chỉ trả lời tốt khi dữ liệu đầu vào đầy đủ, đúng quyền và được cập nhật định kỳ. Các nhóm dữ liệu nên chuẩn hóa trước khi triển khai gồm:

\begin{longtable}{p{3.2cm}p{5.2cm}p{5.2cm}}
\toprule
\textbf{Nhóm dữ liệu} & \textbf{Ví dụ} & \textbf{Dùng cho AI nào} \\
\midrule
DMKH & Tên, địa chỉ, kênh bán, nhóm khách, khu vực, tọa độ & Gợi ý tuyến, scoring, phân quyền khách hàng \\
Lịch sử bán hàng & Ngày bán, sản phẩm, số lượng, doanh thu, tần suất mua & Gợi ý đơn hàng, dự đoán chu kỳ mua \\
Sản phẩm & Mã hàng, tên hàng, nhóm hàng, hoạt chất, đơn vị tính, giá & Tra cứu sản phẩm, catalog, upsell \\
Tồn kho & Kho chính, DL02, DL03, lô, hạn dùng, số lượng khả dụng & Tra tồn kho, đề xuất xả hàng, cảnh báo cận date \\
Tuyến sale & TDV phụ trách, ngày ghé, tần suất, lịch sử check-in & Gợi ý khách nên ghé, giám sát tuyến \\
CTBH/CTKM & Chiết khấu, thưởng, tích lũy, thời gian áp dụng & Gợi ý bán thêm, nhắc mốc thưởng \\
Catalog & Công dụng, hình ảnh, quy cách, hướng dẫn tư vấn & AI catalog, tư vấn nhanh tại điểm bán \\
Công nợ & Nợ tổng, nợ chi tiết, hạn thanh toán, hóa đơn liên quan & AI công nợ, nhắc thu hồi nợ \\
Tài liệu nội bộ & PDF, Excel, chính sách, quy trình, bảng câu hỏi & AI hỏi đáp chính sách và đào tạo \\
\bottomrule
\end{longtable}

\begin{tcolorbox}[warnBox]
Dữ liệu nhạy cảm như công nợ, doanh số và tồn kho phải đi qua lớp phân quyền. TDV chỉ xem dữ liệu cá nhân hoặc khách phụ trách; Manager xem dữ liệu nhóm; Admin/CEO xem theo phạm vi được cấp.
\end{tcolorbox}

% ============================================================================
\part{Hướng dẫn theo vai trò}
% ============================================================================

\chapter{Dành cho TDV / Sale}
\screenuser{Trình dược viên, nhân viên kinh doanh, nhân viên chăm sóc điểm bán.}
\screenpurpose{Giúp TDV hoàn thành công việc trong ngày: xem tuyến, ghé khách, hỏi AI, tra tồn kho, lên đơn và theo dõi công nợ.}

\section{Công việc chính}
\begin{enumerate}
  \item Đăng nhập và xem Dashboard cá nhân.
  \item Xem tuyến hôm nay và danh sách khách nên ghé.
  \item Check-in tại nhà thuốc hoặc điểm bán.
  \item Hỏi AI nên bán sản phẩm nào cho từng khách.
  \item Tra cứu tồn kho, catalog, CTBH và công nợ.
  \item Tạo đơn hàng, theo dõi đơn đã tạo và doanh thu cá nhân.
\end{enumerate}

\section{Flow một ngày của TDV}
\begin{tcolorbox}[roleBox]
\textbf{Buổi sáng:} Đăng nhập $\rightarrow$ xem Dashboard $\rightarrow$ xem Top khách nên ghé $\rightarrow$ xem CTBH đang chạy.

\textbf{Khi đi tuyến:} Mở Tuyến $\rightarrow$ xem khách sắp hết hàng $\rightarrow$ check-in $\rightarrow$ hỏi AI nên bán gì $\rightarrow$ tra tồn kho $\rightarrow$ lên đơn.

\textbf{Cuối ngày:} Xem doanh thu $\rightarrow$ xem đơn đã tạo $\rightarrow$ xem khách chưa chăm sóc $\rightarrow$ xem công nợ cần nhắc.
\end{tcolorbox}

\section{Câu hỏi AI mẫu cho TDV}
\prompt{Hôm nay tôi nên ghé khách nào trước?}
\prompt{Khách Nhà thuốc Minh Châu nên bán thêm sản phẩm nào?}
\prompt{Kiểm tra tồn kho Gabacitin ở kho chính và DL02.}
\prompt{Khách này còn thiếu bao nhiêu để đạt mốc thưởng tháng này?}
\prompt{Tạo gợi ý đơn hàng cho khách ABC dựa trên lịch sử mua 6 tháng gần nhất.}

\chapter{Dành cho Manager / Quản lý vùng}
\screenuser{Quản lý vùng, quản lý nhóm TDV, trưởng bộ phận kinh doanh.}
\screenpurpose{Giúp Manager giám sát hiệu suất nhóm, phát hiện khách giảm mua, duyệt đơn và theo dõi CTBH.}

\section{Công việc chính}
\begin{enumerate}
  \item Xem Dashboard toàn vùng hoặc toàn nhóm.
  \item Theo dõi doanh số, KPI, check-in và hoạt động TDV.
  \item Xem khách lâu chưa phát sinh đơn, khách nguy cơ giảm doanh số và khách có khả năng tăng trưởng.
  \item Duyệt đơn hàng, duyệt phiếu trả và theo dõi trạng thái xử lý.
  \item Theo dõi CTBH, sản phẩm trọng tâm và danh sách khách sắp đạt thưởng.
\end{enumerate}

\section{Bảng tra cứu nhanh cho Manager}
\begin{longtable}{p{6cm}p{8cm}}
\toprule
\textbf{Manager muốn biết} & \textbf{Vào đâu hoặc hỏi AI gì} \\
\midrule
TDV nào dưới KPI & Dashboard quản lý / Báo cáo doanh số theo TDV \\
Khách nào 45 ngày không mua & AI cảnh báo khách giảm mua / Customer Scoring \\
Khách nào sắp đạt thưởng & AI CTBH / Báo cáo chương trình bán hàng \\
Khách nào nguy cơ rời bỏ & Customer Scoring / RFM-C \\
TDV nào chưa ghé tuyến & Giám sát check-in / Lịch sử tuyến \\
Sản phẩm nào đang bán chậm & Báo cáo sản phẩm / AI tồn kho \\
\bottomrule
\end{longtable}

\section{Câu hỏi AI mẫu cho Manager}
\prompt{Liệt kê các khách hàng 45 ngày chưa phát sinh đơn trong vùng của tôi.}
\prompt{TDV nào đang thấp hơn KPI tháng này?}
\prompt{Khách nào có khả năng đạt thưởng nếu mua thêm trong tuần này?}
\prompt{Tổng hợp doanh số theo TDV và nhóm sản phẩm trọng tâm.}

\chapter{Dành cho Admin}
\screenuser{Admin nghiệp vụ, điều phối kinh doanh, người quản trị dữ liệu.}
\screenpurpose{Giúp Admin cấu hình dữ liệu, người dùng, phân quyền, CTBH và tri thức AI để hệ thống vận hành đúng.}

\section{Công việc chính}
\begin{enumerate}
  \item Quản lý người dùng và phân quyền theo vai trò, vùng, kho, nhóm khách.
  \item Khai báo kho chính, DL02, DL03 và phạm vi xem tồn kho.
  \item Import danh mục khách hàng, sản phẩm, tuyến bán hàng và catalog.
  \item Khai báo CTBH, CTKM, chương trình tích lũy, thưởng, chiết khấu.
  \item Upload PDF/Excel/chính sách nội bộ vào kho tri thức AI.
  \item Quản lý ngân hàng câu hỏi AI và thời gian áp dụng từ ngày đến ngày.
\end{enumerate}

\section{Checklist trước khi mở hệ thống cho người dùng}
\begin{itemize}
  \item \checkbox Tài khoản người dùng đã có đúng vai trò.
  \item \checkbox TDV đã được gán Manager, khu vực và tuyến bán hàng.
  \item \checkbox Kho và phân quyền tồn kho đã đúng phạm vi.
  \item \checkbox CTBH/CTKM hiện hành đã có ngày bắt đầu và ngày kết thúc.
  \item \checkbox Catalog sản phẩm có đủ tên, mã, nhóm hàng và mô tả tư vấn.
  \item \checkbox Dữ liệu mẫu đã được kiểm tra trên Dashboard và AI Copilot.
\end{itemize}

\chapter{Dành cho CEO / Ban lãnh đạo}
\screenuser{CEO, Ban giám đốc, Giám đốc kinh doanh, lãnh đạo vận hành.}
\screenpurpose{Giúp lãnh đạo nắm tình hình kinh doanh, tồn kho và cơ hội tăng trưởng bằng dữ liệu tổng hợp.}

\section{Công việc chính}
\begin{enumerate}
  \item Xem tổng quan doanh số theo vùng, TDV, nhóm sản phẩm và thời gian.
  \item Xem vòng quay sản phẩm, tồn kho chậm luân chuyển và sản phẩm cận date.
  \item Yêu cầu AI đề xuất CTKM, combo bán hàng hoặc kế hoạch xả hàng.
  \item Theo dõi khách hàng tăng trưởng, suy giảm hoặc có nguy cơ rời bỏ.
  \item Phát hiện vùng, TDV hoặc sản phẩm có vấn đề để giao việc cho Manager.
\end{enumerate}

\section{Prompt điều hành mẫu}
\prompt{Sản phẩm nào nên chạy combo tháng này?}
\prompt{Sản phẩm nào nên xả hàng vì cận date hoặc tồn kho cao?}
\prompt{Nhóm khách nào đang giảm mua mạnh trong 3 tháng gần đây?}
\prompt{Đề xuất chương trình khuyến mãi cho Gabacitin tháng này.}

\chapter{Dành cho IT / System Admin}
\screenuser{IT nội bộ, đội triển khai, đội tích hợp, người vận hành hệ thống.}
\screenpurpose{Giúp IT đảm bảo hệ thống kết nối đúng dữ liệu, phân quyền đúng và AI hoạt động ổn định.}

\section{Phạm vi trách nhiệm}
\begin{enumerate}
  \item Cấu hình kết nối database, API gateway, n8n và các endpoint tích hợp.
  \item Kiểm tra mapping user, EmployeeID, ManagerID, BranchID và quyền kho.
  \item Theo dõi log lỗi, timeout, retry và chất lượng phản hồi AI.
  \item Bảo vệ secret, chuỗi kết nối, token và quyền truy cập hệ thống.
  \item Sao lưu dữ liệu cấu hình, tài liệu AI và báo cáo định kỳ.
\end{enumerate}

% ============================================================================
\part{Module chức năng chính}
% ============================================================================

\chapter{Module CRM nghiệp vụ}

\section{Dashboard}
\modulecard{Dashboard cá nhân và Dashboard quản lý}{TDV, Manager, CEO}{Theo dõi doanh số, đơn hàng, KPI, tuyến, cảnh báo và tiến độ CTBH trên một màn hình tổng quan.}

\section{Tuyến bán hàng và check-in}
\modulecard{Quản lý tuyến bán hàng}{TDV, Manager}{Hiển thị khách cần ghé trong ngày, ưu tiên theo dữ liệu AI, hỗ trợ check-in và giám sát lịch sử chăm sóc.}

\section{Đơn hàng và phiếu trả}
\modulecard{Quản lý đơn hàng}{TDV, Manager}{Tạo đơn, theo dõi trạng thái, duyệt đơn, ghi nhận phiếu trả và kiểm soát quy trình xử lý.}

\section{Khách hàng và scoring}
\modulecard{Quản lý khách hàng}{TDV, Manager}{Lưu thông tin nhà thuốc, phân nhóm khách, đánh giá RFM-C, phát hiện khách VIP, ổn định hoặc nguy cơ.}

\section{Doanh số, công nợ và hóa đơn}
\modulecard{Báo cáo bán hàng}{TDV, Manager, CEO}{Tra cứu doanh thu, hóa đơn, công nợ, lịch sử mua và tiến độ hoàn thành chỉ tiêu.}

\chapter{Module AI Copilot}

\section{Module 1 — AI Gợi ý đơn hàng}
\screenpurpose{Giúp TDV biết hôm nay nên bán sản phẩm nào cho từng khách hàng.}
\textbf{AI dựa vào:} lịch sử mua, chu kỳ mua, mùa vụ, CTKM hiện tại, sản phẩm trọng tâm, tồn kho và hành vi khách tương tự.

\begin{tcolorbox}[exBox]
\textbf{Khách hàng:} Nhà thuốc Minh Châu

\textbf{AI đề xuất:}
\begin{enumerate}
  \item Gabacitin
  \item Bảo Can Linh
  \item Combo Q10 + Bổ thần kinh
\end{enumerate}

\textbf{Lý do:} Khách từng mua Gabacitin, chu kỳ mua trung bình 2--3 tháng, sản phẩm đang nằm trong CTBH trọng tâm và tồn kho còn đủ để bán.
\end{tcolorbox}

\section{Module 2 — AI Tuyến bán hàng}
AI sắp xếp khách nên ghé theo mức ưu tiên: khách lâu chưa mua, khách sắp hết hàng, khách gần vị trí TDV, khách có khả năng đạt thưởng, khách có công nợ cần nhắc hoặc khách đang trong CTBH.

\prompt{Hôm nay tôi nên đi tuyến nào để tối ưu doanh số?}
\prompt{Sắp xếp 10 khách cần ghé theo mức ưu tiên trong khu vực của tôi.}

\section{Module 3 — Customer Scoring RFM-C}
RFM-C đánh giá khách hàng theo bốn nhóm tín hiệu:
\begin{itemize}
  \item \textbf{Recency}: lần mua gần nhất.
  \item \textbf{Frequency}: tần suất mua.
  \item \textbf{Monetary}: doanh số hoặc giá trị mua.
  \item \textbf{Consumption}: sản lượng tiêu thụ và tốc độ dùng hàng.
\end{itemize}

\begin{longtable}{p{3cm}p{5cm}p{6cm}}
\toprule
\textbf{Nhóm} & \textbf{Ý nghĩa} & \textbf{Hành động gợi ý} \\
\midrule
A / VIP & Mua đều, doanh số cao, khả năng tăng trưởng tốt & Ưu tiên chăm sóc, gợi ý combo, nhắc CTBH \\
B / Ổn định & Mua đều nhưng chưa tối ưu giá trị đơn & Gợi ý upsell, theo dõi chu kỳ mua \\
C / Nguy cơ & Lâu chưa mua, giảm doanh số hoặc giảm tần suất & Manager/TDV cần gọi lại, ghé thăm hoặc áp dụng CTKM \\
\bottomrule
\end{longtable}

\section{Module 4 — Tích lũy, thưởng và CTBH}
AI theo dõi tiến độ khách hàng trong từng chương trình bán hàng: đã mua bao nhiêu, còn thiếu bao nhiêu, cần mua thêm sản phẩm nào để đạt mốc thưởng hoặc chiết khấu.

\prompt{Khách ABC còn thiếu bao nhiêu để đạt mốc thưởng tháng này?}
\prompt{Gợi ý sản phẩm bán thêm để khách đạt chiết khấu tốt hơn.}

\section{Module 5 — Upsell theo triệu chứng và catalog}
AI hỗ trợ TDV tra cứu sản phẩm theo triệu chứng, nhóm bệnh, công dụng, hoạt chất hoặc nhu cầu tư vấn của nhà thuốc.

\prompt{Khách hỏi sản phẩm hỗ trợ mất ngủ, nên tư vấn sản phẩm nào?}
\prompt{Tìm sản phẩm thay thế cho nhóm bổ thần kinh đang có tồn kho tốt.}

\section{Module 6 — Đề xuất khuyến mãi và xả hàng}
Dành cho Manager/CEO/Admin. AI phân tích tồn kho lớn, hàng cận date, sản phẩm bán chậm, vòng quay hàng và biên lợi nhuận để đề xuất CTKM hoặc combo.

\prompt{Đề xuất combo xả hàng cho nhóm sản phẩm cận date trong 60 ngày tới.}
\prompt{Sản phẩm nào tồn kho cao nhưng bán chậm trong 3 tháng gần đây?}

\section{Module 7 — Công nợ và hóa đơn}
AI hỗ trợ tra cứu công nợ tổng, công nợ chi tiết, hóa đơn liên quan và danh sách khách cần nhắc thanh toán.

\prompt{Khách Nhà thuốc Minh Châu còn nợ bao nhiêu?}
\prompt{Liệt kê các khách có công nợ quá hạn trong tuyến của tôi.}

\section{Module 8 — Gợi ý đơn thuốc và xử lý ảnh}
Khi TDV chụp toa thuốc hoặc hóa đơn, AI có thể trích xuất nội dung, nhận diện sản phẩm liên quan và gợi ý sản phẩm thay thế hoặc bán kèm theo danh mục Medstand.

\begin{tcolorbox}[warnBox]
Kết quả nhận diện ảnh chỉ dùng để hỗ trợ thao tác bán hàng. Người dùng cần kiểm tra lại tên sản phẩm, hàm lượng và số lượng trước khi tạo đơn.
\end{tcolorbox}

\section{Module 9 — AI hỏi đáp chính sách nội bộ}
AI đọc tài liệu PDF/Excel/chính sách đã được Admin upload để trả lời câu hỏi về quy trình, CTBH, điều kiện chiết khấu, phân quyền và hướng dẫn vận hành.

\section{Module 10 — Sản phẩm trọng tâm tháng}
Module này theo dõi danh sách sản phẩm trọng tâm do BGĐ/Admin khai báo, tiến độ bán theo TDV/khách hàng và cơ hội tăng doanh số.

\section{Module 11 — Tra cứu sản phẩm và tồn kho}
AI trả lời tồn kho theo kho được phân quyền, gồm kho chính, DL02, DL03, lô hàng và hạn dùng nếu dữ liệu có sẵn.

\prompt{Tồn kho Gabacitin ở kho chính, DL02 và DL03 hiện còn bao nhiêu?}

\section{Module 12 — Khảo sát khách hàng}
Hệ thống hỗ trợ tạo bài khảo sát, ghi nhận câu trả lời và tổng hợp kết quả để Manager/Admin đánh giá thị trường, chất lượng chăm sóc hoặc hiểu biết sản phẩm của TDV.

% ============================================================================
\part{Vận hành, triển khai và kiểm thử}
% ============================================================================

\chapter{Quy trình triển khai dữ liệu}

\section{Bước 1 — Chuẩn hóa dữ liệu nền}
\begin{enumerate}
  \item Chuẩn hóa mã khách hàng, mã sản phẩm, mã nhân viên và mã kho.
  \item Kiểm tra trùng lặp khách hàng và sai khác tên nhà thuốc.
  \item Đảm bảo lịch sử bán hàng tối thiểu 6 tháng, khuyến nghị 12 tháng.
  \item Gán TDV, Manager, vùng, tuyến và quyền kho.
\end{enumerate}

\section{Bước 2 — Nạp dữ liệu nghiệp vụ}
\begin{enumerate}
  \item Nạp DMKH, danh mục sản phẩm, tồn kho, hóa đơn, đơn hàng.
  \item Nạp CTBH, CTKM, catalog và chính sách nội bộ.
  \item Kiểm tra dashboard, báo cáo, AI Copilot và phân quyền mẫu.
\end{enumerate}

\section{Bước 3 — Kiểm thử theo vai trò}
\begin{longtable}{p{3cm}p{5cm}p{6cm}}
\toprule
\textbf{Vai trò} & \textbf{Cần kiểm thử} & \textbf{Kết quả mong đợi} \\
\midrule
TDV & Login, tuyến, đơn hàng, tồn kho, công nợ, AI gợi ý & Chỉ thấy dữ liệu cá nhân và khách phụ trách \\
Manager & Dashboard vùng, dữ liệu TDV dưới quyền, duyệt đơn & Thấy dữ liệu nhóm, không vượt quyền ngoài vùng \\
Admin & Import, phân quyền, CTBH, upload tài liệu & Cấu hình có hiệu lực và AI đọc đúng tài liệu \\
CEO & Báo cáo tổng, phân tích tồn kho, đề xuất CTKM & Thấy dữ liệu tổng hợp theo quyền lãnh đạo \\
IT & Kết nối, log, API, bảo mật, backup & Hệ thống ổn định, không lộ secret, log đủ để truy vết \\
\bottomrule
\end{longtable}

\chapter{Checklist đào tạo người dùng}

\section{Buổi đào tạo TDV}
\begin{itemize}
  \item \checkbox Đăng nhập và đổi thông tin cá nhân.
  \item \checkbox Xem Dashboard và hiểu các chỉ số chính.
  \item \checkbox Mở tuyến, xem khách nên ghé, check-in.
  \item \checkbox Hỏi AI gợi ý đơn hàng và tra tồn kho.
  \item \checkbox Tạo đơn hàng mẫu và kiểm tra trạng thái.
  \item \checkbox Tra công nợ và lịch sử mua của khách.
\end{itemize}

\section{Buổi đào tạo Manager}
\begin{itemize}
  \item \checkbox Xem Dashboard nhóm/vùng.
  \item \checkbox Theo dõi KPI, doanh số, check-in TDV.
  \item \checkbox Duyệt đơn hàng và phiếu trả.
  \item \checkbox Hỏi AI danh sách khách giảm mua, khách sắp đạt thưởng.
  \item \checkbox Đọc báo cáo sản phẩm trọng tâm và CTBH.
\end{itemize}

\section{Buổi đào tạo Admin/IT}
\begin{itemize}
  \item \checkbox Quản lý user, vai trò, vùng, kho.
  \item \checkbox Import dữ liệu và kiểm tra lỗi import.
  \item \checkbox Cấu hình CTBH, CTKM, thời gian áp dụng.
  \item \checkbox Upload tài liệu vào AI Knowledge Base.
  \item \checkbox Kiểm tra log, kết nối và quyền truy cập.
\end{itemize}

% ============================================================================
\appendix
\part{Phụ lục}
% ============================================================================

\chapter{Bảng thuật ngữ}
\begin{longtable}{p{4cm}p{10cm}}
\toprule
\textbf{Thuật ngữ} & \textbf{Giải thích} \\
\midrule
TDV & Trình dược viên / nhân viên kinh doanh phụ trách nhà thuốc hoặc điểm bán. \\
Manager & Quản lý vùng hoặc quản lý nhóm TDV. \\
CTBH & Chương trình bán hàng, gồm thưởng, chiết khấu, tích lũy hoặc sản phẩm trọng tâm. \\
CTKM & Chương trình khuyến mãi. \\
RFM-C & Mô hình chấm điểm khách hàng theo Recency, Frequency, Monetary và Consumption. \\
AI Copilot & Trợ lý AI hỗ trợ hỏi đáp, gợi ý, phân tích và tra cứu dữ liệu. \\
Knowledge Base & Kho tri thức gồm tài liệu, chính sách, PDF, Excel hoặc nội dung nội bộ để AI tra cứu. \\
DL02/DL03 & Các kho/điểm lưu kho hoặc đại lý được dùng trong phân quyền tồn kho. \\
\bottomrule
\end{longtable}

\chapter{FAQ}

\section*{Dành cho TDV}
\textbf{Tôi không thấy khách trong tuyến hôm nay thì làm sao?} Kiểm tra bộ lọc ngày/khu vực trước. Nếu vẫn thiếu, báo Manager hoặc Admin kiểm tra phân tuyến.

\textbf{AI gợi ý sản phẩm nhưng tồn kho không đủ thì sao?} Ưu tiên kiểm tra tồn kho trước khi lên đơn. Có thể hỏi AI sản phẩm thay thế hoặc kho khác còn hàng.

\textbf{AI trả lời sai tên khách hoặc sai dữ liệu thì sao?} Không tạo đơn ngay. Chụp màn hình, ghi lại câu hỏi đã dùng và gửi Admin/IT kiểm tra dữ liệu nguồn.

\section*{Dành cho Manager}
\textbf{Tôi không thấy dữ liệu TDV dưới quyền?} Kiểm tra mapping ManagerID, EmployeeID, BranchID và trạng thái tài khoản TDV.

\textbf{Số liệu Dashboard lệch với báo cáo kế toán?} Kiểm tra khoảng thời gian, trạng thái đơn/hóa đơn, nguồn dữ liệu và thời điểm đồng bộ.

\section*{Dành cho Admin/IT}
\textbf{Upload tài liệu nhưng AI chưa trả lời được?} Kiểm tra định dạng file, trạng thái index, quyền truy cập tài liệu và thời gian đồng bộ tri thức.

\textbf{Người dùng thấy dữ liệu vượt quyền?} Tạm khóa quyền hoặc role liên quan, kiểm tra điều kiện lọc theo user, ManagerID, BranchID và quyền kho.

\chapter{Xử lý sự cố}

\section*{Không đăng nhập được}
\begin{enumerate}
  \item Kiểm tra đúng username/password.
  \item Kiểm tra tài khoản có bị Disable hay không.
  \item Kiểm tra kết nối mạng và API đăng nhập.
  \item Nếu vẫn lỗi, gửi ảnh màn hình và thời điểm lỗi cho IT.
\end{enumerate}

\section*{Không thấy dữ liệu}
\begin{enumerate}
  \item Kiểm tra bộ lọc ngày, vùng, kho, TDV, trạng thái đơn hàng.
  \item Kiểm tra user đã được gán EmployeeID, ManagerID, BranchID chưa.
  \item Kiểm tra dữ liệu đã đồng bộ từ ERP/database sang hệ thống chưa.
\end{enumerate}

\section*{AI phản hồi chậm hoặc timeout}
\begin{enumerate}
  \item Rút gọn câu hỏi, nêu rõ khách hàng/sản phẩm/thời gian.
  \item Thử lại sau vài phút nếu hệ thống đang tải cao.
  \item Gửi log câu hỏi, thời điểm và màn hình đang thao tác cho IT.
\end{enumerate}

\chapter{Kịch bản demo nhanh}

\section*{Demo 15 phút cho khách hàng}
\begin{enumerate}
  \item Mở Dashboard để giới thiệu tổng quan doanh số và KPI.
  \item Mở Tuyến bán hàng và cho AI gợi ý khách nên ghé.
  \item Chọn một khách hàng, hỏi AI nên bán sản phẩm nào.
  \item Tra tồn kho sản phẩm được gợi ý.
  \item Tạo đơn hàng mẫu.
  \item Chuyển sang Manager để xem dữ liệu nhóm và duyệt đơn.
  \item Hỏi AI sản phẩm nào nên chạy CTKM hoặc xả hàng.
\end{enumerate}

\begin{tcolorbox}[tipBox]
Một demo tốt nên đi theo câu chuyện kinh doanh: \textbf{khách nào cần ghé} $\rightarrow$ \textbf{bán gì} $\rightarrow$ \textbf{còn hàng không} $\rightarrow$ \textbf{lên đơn} $\rightarrow$ \textbf{Manager theo dõi được gì}.
\end{tcolorbox}

\chapter*{Kết thúc tài liệu}
\addcontentsline{toc}{chapter}{Kết thúc tài liệu}

Tài liệu v3.0 này là nền để tiếp tục bổ sung ảnh chụp màn hình, quy trình thao tác chi tiết theo từng màn hình và kịch bản đào tạo thực tế sau khi giao diện hệ thống được chốt.

\printindex

\end{document}

